% =====================================================================
% Paper 2, full article (pre-registered study), Nature-family Article format:
% main text <= 3,500 words, abstract <= 150 words, <= 6 display items, ~50 references
% (nature.com/natcomputsci and /natchemeng content types, read live 2026-10-20).
% Status: METHODS drafted 2026-11-01 while the A6 rerun ran; Introduction, Results and
% Discussion are written only after the corrected verdicts exist. No result appears below.
% Source of truth: PREREG_P2_discoverability.md (frozen 7dcc2fe; amendments A1-A6 in its s.9).
% Number provenance: every number is quoted from the prereg, a results/test/log file or the
% code, named in a comment beside it. Writing rules: plain prose, at most five em dashes in
% the whole paper, no claim beyond the data.
% =====================================================================
\documentclass[11pt]{article}
\usepackage[margin=2.5cm]{geometry}
\usepackage[T1]{fontenc}
\usepackage[utf8]{inputenc}
\usepackage{lmodern}
\usepackage{amsmath,amssymb}
\usepackage{booktabs}
\usepackage{array}
\usepackage{tabularx}
\usepackage{graphicx}
\usepackage{setspace}
\usepackage{lineno}
\usepackage[numbers,sort&compress]{natbib}
\usepackage{xcolor}
\usepackage[colorlinks=true,linkcolor=blue!50!black,citecolor=blue!50!black,urlcolor=blue!50!black]{hyperref}
\onehalfspacing
\linenumbers

\newcommand{\Da}{\mathrm{Da}}
\newcommand{\Pe}{\mathrm{Pe}}
\newcommand{\qs}{q^{*}}
\newcommand{\qsm}{q^{*}_{\mathrm{meas}}}
\newcommand{\verdict}[1]{``\textit{#1}''}

\title{When can AI discover adsorption kinetics?}
\author{Abdulhamid Batayhi\\[4pt]
\small Department of Mechatronics Engineering, Faculty of Mechanical Engineering,\\
\small Y\i{}ld\i{}z Technical University, Istanbul, T\"urkiye\\[4pt]
\small ORCID 0009-0007-0374-6755}
\date{}

\begin{document}
\maketitle

% Abstract, Introduction, Results and Discussion: written after the corrected verdicts.

\section*{Methods}

\subsection*{Ground truth}
Ground truth came from the finite-difference reference solver of the companion
study \citep{paper1}: a one-dimensional, axially dispersed, non-isothermal model of a
packed column, verified there against closed-form solutions. Uptake followed one of two
laws,
\begin{align}
\text{L1:}\quad & \frac{\partial q}{\partial t} = k\,\bigl(\qs(c,T) - q\bigr),\\
\text{L2:}\quad & \frac{\partial q}{\partial t} = k(q)\,\bigl(\qs(c,T) - q\bigr),\qquad
k(q) = k_0\Bigl[1 - 0.8\,s\!\Bigl(\frac{q/q_{\max}-\theta}{0.05}\Bigr)\Bigr],
\end{align}
where $s$ is the logistic function and $\theta$ is the fractional loading at the
isotherm's step midpoint, so that under L2 the rate falls to 20\,\% of $k_0$ across the step.
% PREREG_P2 s.4.1; p2_generate.py (L2_DROP 0.8, L2_WIDTH 0.05, step_midpoint_fraction).
Two adsorbents were simulated with the same column geometry and flow (Supplementary
Table~1): the companion study's favourable Langmuir sorbent, and water on a MOF-303-like
framework, a cooperative (Type V) isotherm with its step at $P/P_0 = 0.15$, fed at 30\,\%
relative humidity and 30\,$^\circ$C. The Damk\"ohler number $\Da = k\,t_{\mathrm{stoich}}$
counts kinetic time constants per passage of the adsorption front, with
$t_{\mathrm{stoich}}$ the stoichiometric time of the column (7200\,s for Langmuir and
31\,050\,s for the step isotherm). Nine levels from $10^{-1}$ to $10^{3}$ at half-decade
spacing therefore span $k$ from $1.39\times10^{-5}$ to $0.139$\,s$^{-1}$ for Langmuir and
from $3.22\times10^{-6}$ to $0.0322$\,s$^{-1}$ for the step isotherm. Three axial-dispersion
coefficients, $D_L \times 0.2$, $\times 1$ and $\times 5$ of the default, completed the grid.
They give column P\'eclet numbers $\Pe = uL/D_L$ of about 470, 94 and 19, and the levels are
named by these values below (the code and the pre-registration label them by the $D_L$
multiplier, so its ``$\times 0.2$'' level is the highest $\Pe$). The grid has 81 solves at 400 cells and 400 stored times, each run until the outlet reached 95\,\%
breakthrough, plus 9 isothermal solves for a control described below. The manifest that
records every parameter is public.
% data/p2/manifest.json (t_stoich_s 7200.4 and 31050.4; k_LDF 1.3888e-05..0.13888 and
% 3.2206e-06..0.032206); p2_generate.py N_Z = 400, N_SNAP = 400, 95 % breakthrough rule;
% p2_generate.log "done: 81 cells; without 95 % breakthrough: none".

\subsection*{Observation models, noise and isotherm error}
Every method and every identifiability calculation read the data through one function,
so all saw the same corrupted observations. Three observation models were used: interior
probes of $c$, $q$ and $T$ (O1), interior probes of $c$ and $T$ with $q$ hidden (O2), and the
outlet $c(t)$ and $T(t)$ only (O3). O1 and O2 used 20 equally spaced probes; all models
used 200 times spread evenly over the solved horizon. Noise was independent Gaussian with
standard deviation $\sigma \in \{0, 0.5, 2, 5\}\,\%$ of each channel's range. Methods were
never given the true isotherm. They were given a measured one,
$\qsm(c,T) = \qs(c,T)\,[1 + \varepsilon\, g(c/c_{\mathrm{in}})]$ with
$\varepsilon \in \{0, 0.5, 2, 5\}\,\%$, where $g$ is a smooth random function of four
Fourier modes on $[0, 1]$, scaled to unit root-mean-square and drawn once per replicate.
Every random draw was seeded by a hash of (cell, replicate, $\sigma$, $\varepsilon$,
observation model), so each observation can be regenerated on its own. A replicate, one
such observation of one cell, was the unit of analysis; each cell had 60.
% p2_observe.py (N_PROBE 20, N_T 200, N_MODES 4, SHA-256 seeding, isotherm_error_fn);
% PREREG_P2 s.4.2 and A1 (60 replicates).

\subsection*{Discovery methods}
Each method ran in the strongest configuration we could verify on known-answer systems
before the pre-registration was frozen, and a method that could not be run that way was
reported as such rather than replaced by a weaker one. The candidate library was
$\{1, c, q, T, cq, q^2, c^2, \qsm\}$; under L1 the true law is $k(\qsm - q)$. The methods
were sparse regression on a Savitzky--Golay estimate of $\mathrm{d}q/\mathrm{d}t$ (strong
form) and on its weak, integral form with compactly supported test functions
\citep{messenger2021ode,messenger2021pde}, each with three solvers: sequentially thresholded
least squares \citep{brunton2016}, a bagged ensemble with an inclusion threshold of 0.6
\citep{fasel2022}, and exact best-subset selection by enumeration with the Bayesian
information criterion (M4, the objective of \citep{bertsimas2022}); errors-in-variables
best-subset regression given the true error levels (M5, an oracle); implicit sparse
regression without any isotherm on an isothermal control (M6, SINDy-PI \citep{kaheman2020});
symbolic regression (M7, PySR \citep{cranmer2023}); a universal differential equation with a
neural network in the uptake term, trained through the column model and read out by sparse
and symbolic regression (M1, the pipeline of \citep{santana2023}); and regression of the state
on $q$-free features, which returns the isotherm where uptake is at equilibrium (M9
\citep{delgado2025}). On O2, uptake was first recovered from the gas-phase mass balance,
and every O1 method was then run on it. M1 could be trained for at most 50 L-BFGS steps or
2\,h per fit on the machine available, so it ran as a declared demonstration: Langmuir,
$\Pe \approx 94$, $\sigma = 2\,\%$, $\varepsilon = 0$, O3, the nine $\Da$ levels, one
replicate. The best parameters seen during training were kept. As first run, a column solve
that exhausted its step budget during training ended the fit with no readout (four of the
nine fits); the corrected code keeps the best parameters in that case, and both runs are
reported (Supplementary Note~2). A Kolmogorov--Arnold network \citep{liu2024} failed its known-answer test and its
cost gate before the freeze and was not run on the grid. Method settings, solver thresholds
and their known-answer tests are listed in Supplementary Note~1.
% PREREG_P2 s.3 and A4; p2/methods.py, p2/weak.py, p2/eiv.py, p2/pysr_method.py,
% p2_grid_extra.py (fit_m1: M1_MAXITER 50, M1_BUDGET_S 7200; _train_lbfgs, B88);
% results/p2_extra_m1_asrun.json; results/p2_cost_gates.json.

\subsection*{Success}
Discovery succeeded when the selected terms were exactly those of the true law, every sign
was correct and, for L1, the ratio of the $\qsm$ and $q$ coefficients was $-1$ within
10\,\%. The value of $k$ itself was not required. A method that raised an error counted as
a failure; no replicate, cell or output was excluded.
% p2/metric.py success_L1; PREREG_P2 s.4.3 and A2 (errors as failures).

\subsection*{Identifiability}
Practical identifiability of $k$ was judged by profile likelihood \citep{raue2009} with the
true law known, on the same observations and with the same measured isotherm as the
discovery methods: $k$ was practically identifiable when the 95\,\% profile interval
(likelihood-ratio threshold 3.84) lay within $[k/2, 2k]$. On O1 the likelihood summed over
all probes the fit of the known law, integrated exactly for piecewise-linear $\qs(t)$ along
the observed $c$ and $T$. On O2 and O3 every point of the profile was a full column solve.
The profile searched a coarse logarithmic grid of $k$ and refined it around the interval.
Two corrections were made during the study and are reported in full (Supplementary
Note~2). First, on O2 and O3 a measured isotherm with $\varepsilon > 0$ could stall the stiff
solver at large trial $k$ for days; each solve was therefore capped at 600\,s and a capped
solve was scored as a failed solve (amendment A5; two O3 rows had capped solves, none next to
its interval). Second, the O1 grid first spanned $10^{-5}$ to $10$\,s$^{-1}$, which excluded the
step cells at $\Da = 0.1$; it was widened to $10^{-7}$ to $10$\,s$^{-1}$, a peak on the grid
edge was made an error, and all O1 rows were recomputed (amendment A6).
% PREREG_P2 s.4.3; p2/identifiability.py (profile_interval_o1, O1_K_GRID, _profile,
% SOLVE cap); p2_grid_o23.py SOLVE_CAP_S 600; results/p2_o3_ident.json (n_capped);
% RETRACTIONS_P2 B86, B87.

\subsection*{Boundaries and verdicts}
For each slice the boundary $\Da^{*}$ was the 50\,\% crossing of an unpenalised logistic
regression of replicate outcomes on $\log_{10}\Da$, with a percentile bootstrap interval
(2000 resamples within each $\Da$ level). No boundary was reported, and no verdict issued,
when the outcomes did not cross 50\,\% inside the tested range; the range was reported
instead. Before the freeze the estimator's coverage on planted boundaries was 0.94 for a
nominal 0.95, with a bias of $+0.012$ decade. The best discovery method of each slice was
chosen on replicates 0--14 and its boundary estimated on replicates 15--59, which removes the
inflation that choosing and estimating on the same replicates causes (0.07 decade in
simulation); the identifiability boundary was estimated on the same 45 replicates. M5, which
knows the true error levels, took no part in this choice.
% p2/analysis.py boundary (n_boot 2000), select_best; p2_analyze.py N_SEL = 15;
% tests/test_p2.py::test_boundary_ci_has_nominal_coverage_and_no_bias (0.94, +0.012);
% results/p2_operating_characteristics.json selection bias_same 0.0708; PREREG_P2 A1, A2.

Three hypotheses were confirmatory, each with verdict words fixed in advance. H2a: success
collapses onto one curve in $R = \delta/\varepsilon_{\mathrm{eff}}$, the driving force
relative to the combined error, with the spread of $R^{*}$ across conditions under half a
decade. H2b: on O1, for each slice ($\sigma$, $\varepsilon$, isotherm, $\Pe$; 96 in all),
the gap $\log_{10}\Da^{*}_{\mathrm{disc}} - \log_{10}\Da^{*}_{\mathrm{ident}}$ was called
\verdict{discovery fails before identifiability} at $-0.25$ decade or less,
\verdict{discovery outlives identifiability} at $+0.25$ or more, and
\verdict{discovery fails with identifiability} in between; the hypothesis-level verdict was
the word held by at least two thirds of the slices with a verdict. H2c: errors-in-variables
regression moves the discovery boundary upwards by at least 0.25 decade in at least two thirds
of the slices at a given $\varepsilon > 0$. With 45 estimation replicates the H2b rule
returned a wrong directional verdict with probability 0.001 when there was no gap, and the
correct one with probability 0.965 at a gap of 0.40 decade, the smallest effect of interest.
A positive control guarded every verdict: in each isotherm's easiest cells (lowest $\Da$,
$\sigma = \varepsilon = 0$, O1) the best method had to succeed and $k$ had to be identifiable
in at least 90\,\% of replicates, or that isotherm's verdicts were withheld. O2 and O3 were
reported descriptively, and H2d, whether the loading-dependent rate of L2 can be discovered,
was exploratory.
% PREREG_P2 s.2 as amended by A1, A2; p2/analysis.py disc_vs_ident (0.25), aggregate_h2b,
% positive_control (need 0.9); results/p2_operating_characteristics.json h2b n45
% (gap 0: 0.001; gap 0.40: 0.965).

\subsection*{Pre-registration and amendments}
The design was frozen on 11 October 2026 (commit \texttt{7dcc2fe}) before any discovery method
ran on any grid cell, and every confirmatory driver refused to run unless the pre-registration
matched its latest freeze commit. Six dated amendments are appended to it. A1 to A4 were made
before the grid ran: 60 replicates split into selection and estimation sets, the smallest
effect of interest, the oracle kept out of the pooled verdicts, $\Pe$ as a slicing axis,
errors counted as failures, the decision to run the grid before peer review, and the M1
demonstration. A5 (the solve cap) was made while the grid ran, and A6 (the O1 grid range)
after the verdicts were first computed. Both as-run and corrected verdicts are public, and
the paper reports both (Supplementary Note~2).
% PREREG_P2 s.9 (A1 ef3ade8, A2 397d9ea, A3 d892ab9, A4 ffc8d88, A5 97150f6, A6 375552b);
% p2_grid_o1.py frozen(); results/p2_verdicts_asrun.json.

\subsection*{Exploratory analyses}
Three analyses were declared in writing before they ran and are reported as exploratory:
M1 trained at the scale a universal differential equation needs (E1), the Kolmogorov--Arnold
network on its declared subgrid once the machine was free (E2), and how often an interval
that meets $[k/2, 2k]$ misses the true $k$ (E3).
% EXPLORATORY_P2.md (committed 787bb29 before any of them ran).

\subsection*{Software and computation}
The grid ran on a single computer without a graphics processor, one job at a time, from
12 October to 1 November 2026. Package versions are pinned in the repository.
% p2_grid_chain.log (start 2026-10-12 17:59:58, P2_GRID_CHAIN_DONE 2026-11-01 03:07:21);
% requirements_p2_venv.txt.

\subsection*{Use of AI tools}
During the preparation of this work the author used a generative AI assistant (Claude,
Anthropic) as a tool to assist with drafting and language editing, with writing and checking
code, and with checking bibliographic records. The author designed the study, made every
scientific decision, reviewed and edited all content, and takes full responsibility for the
content of the published article.

\section*{Data availability}
The ground-truth manifest, every results file, the run log and the as-run and corrected
verdicts are public in the study repository and archived on Zenodo
(doi:10.5281/zenodo.23325994). The ground-truth fields are regenerated from the manifest by
\texttt{p2\_generate.py}.

\section*{Code availability}
All code is public at \url{https://github.com/abdulhamidbatayhi123/ai-kinetics-discoverability}
under the MIT licence and archived with the data. The solver is imported unchanged from the
companion study's archive (doi:10.5281/zenodo.23323685).

\section*{Author contributions}
A.B. is the sole author and conceived and designed the study. The use of AI tools is
described in Methods.

\section*{Competing interests}
The author declares no competing interests.

\begin{thebibliography}{99}
\bibitem{brunton2016} Steven L. Brunton, Joshua L. Proctor, J. Nathan Kutz, Discovering governing equations from data by sparse identification of nonlinear dynamical systems, Proceedings of the National Academy of Sciences, vol. 113, 3932-3937 (2016). \href{https://doi.org/10.1073/pnas.1517384113}{doi:10.1073/pnas.1517384113}; \href{https://arxiv.org/abs/1509.03580}{arXiv:1509.03580}.

\bibitem{rudy2016} Samuel H. Rudy, Steven L. Brunton, Joshua L. Proctor, J. Nathan Kutz, Data-driven discovery of partial differential equations, Science Advances, vol. 3, e1602614 (2017). \href{https://doi.org/10.1126/sciadv.1602614}{doi:10.1126/sciadv.1602614}; \href{https://arxiv.org/abs/1609.06401}{arXiv:1609.06401}.

\bibitem{messenger2021ode} Daniel A. Messenger, David M. Bortz, Weak SINDy: Galerkin-Based Data-Driven Model Selection, Multiscale Modeling \& Simulation, vol. 19, 1474-1497 (2021). \href{https://doi.org/10.1137/20M1343166}{doi:10.1137/20M1343166}; \href{https://arxiv.org/abs/2005.04339}{arXiv:2005.04339}.

\bibitem{messenger2021pde} Daniel A. Messenger, David M. Bortz, Weak SINDy for partial differential equations, Journal of Computational Physics, vol. 443, 110525 (2021). \href{https://doi.org/10.1016/j.jcp.2021.110525}{doi:10.1016/j.jcp.2021.110525}; \href{https://arxiv.org/abs/2007.02848}{arXiv:2007.02848}.

\bibitem{reinbold2020} Patrick A. K. Reinbold, Daniel R. Gurevich, Roman O. Grigoriev, Using noisy or incomplete data to discover models of spatiotemporal dynamics, Physical Review E, vol. 101, 010203 (2020). \href{https://doi.org/10.1103/PhysRevE.101.010203}{doi:10.1103/PhysRevE.101.010203}; \href{https://arxiv.org/abs/1911.03365}{arXiv:1911.03365}.

\bibitem{fasel2022} U. Fasel, J. N. Kutz, B. W. Brunton, S. L. Brunton, Ensemble-SINDy: Robust sparse model discovery in the low-data, high-noise limit, with active learning and control, Proceedings of the Royal Society A: Mathematical, Physical and Engineering Sciences, vol. 478, 20210904 (2022). \href{https://doi.org/10.1098/rspa.2021.0904}{doi:10.1098/rspa.2021.0904}; \href{https://arxiv.org/abs/2111.10992}{arXiv:2111.10992}.

\bibitem{hirsh2021} Seth M. Hirsh, David A. Barajas-Solano, J. Nathan Kutz, Sparsifying priors for Bayesian uncertainty quantification in model discovery, Royal Society Open Science, vol. 9, 211823 (2022). \href{https://doi.org/10.1098/rsos.211823}{doi:10.1098/rsos.211823}; \href{https://arxiv.org/abs/2107.02107}{arXiv:2107.02107}.

\bibitem{kaheman2020} Kadierdan Kaheman, J. Nathan Kutz, Steven L. Brunton, SINDy-PI: a robust algorithm for parallel implicit sparse identification of nonlinear dynamics, Proceedings of the Royal Society A: Mathematical, Physical and Engineering Sciences, vol. 476, 20200279 (2020). \href{https://doi.org/10.1098/rspa.2020.0279}{doi:10.1098/rspa.2020.0279}; \href{https://arxiv.org/abs/2004.02322}{arXiv:2004.02322}.

\bibitem{bertsimas2022} Dimitris Bertsimas, Wes Gurnee, Learning sparse nonlinear dynamics via mixed-integer optimization, Nonlinear Dynamics, vol. 111, 6585-6604 (2023). \href{https://doi.org/10.1007/s11071-022-08178-9}{doi:10.1007/s11071-022-08178-9}; \href{https://arxiv.org/abs/2206.00176}{arXiv:2206.00176}.

\bibitem{cranmer2023} Miles Cranmer, Interpretable Machine Learning for Science with PySR and SymbolicRegression.jl, arXiv preprint arXiv:2305.01582 (2023). \href{https://arxiv.org/abs/2305.01582}{arXiv:2305.01582}.

\bibitem{udrescu2020} Silviu-Marian Udrescu, Max Tegmark, AI Feynman: A physics-inspired method for symbolic regression, Science Advances, vol. 6, eaay2631 (2020). \href{https://doi.org/10.1126/sciadv.aay2631}{doi:10.1126/sciadv.aay2631}; \href{https://arxiv.org/abs/1905.11481}{arXiv:1905.11481}.

\bibitem{santana2023} Vinicius V. Santana, Erbet Costa, Carine M. Rebello, Ana Mafalda Ribeiro, Christopher Rackauckas, Idelfonso B.R. Nogueira, Efficient hybrid modeling and sorption model discovery for non-linear advection-diffusion-sorption systems: A systematic scientific machine learning approach, Chemical Engineering Science, vol. 282, 119223 (2023). \href{https://doi.org/10.1016/j.ces.2023.119223}{doi:10.1016/j.ces.2023.119223}; \href{https://arxiv.org/abs/2303.13555}{arXiv:2303.13555}.

\bibitem{valocchi1985} Albert J. Valocchi, Validity of the Local Equilibrium Assumption for Modeling Sorbing Solute Transport Through Homogeneous Soils, Water Resources Research, vol. 21, 808-820 (1985). \href{https://doi.org/10.1029/wr021i006p00808}{doi:10.1029/wr021i006p00808}.

\bibitem{hansen2017} Scott K. Hansen, Velimir V. Vesselinov, Local equilibrium and retardation revisited, Groundwater, vol. 56, 109-117 (2018). \href{https://doi.org/10.1111/gwat.12566}{doi:10.1111/gwat.12566}; \href{https://arxiv.org/abs/1703.03087}{arXiv:1703.03087}.

\bibitem{katsoulas2023} Konstantinos Katsoulas, Monica Tirapelle, Eva Sorensen, Luca Mazzei, On the apparent dispersion coefficient of the equilibrium dispersion model: An asymptotic analysis, Journal of Chromatography A, vol. 1708, 464345 (2023). \href{https://doi.org/10.1016/j.chroma.2023.464345}{doi:10.1016/j.chroma.2023.464345}.

\bibitem{knox2016} James C. Knox, Armin D. Ebner, M. Douglas LeVan, Robert F. Coker, James A. Ritter, Limitations of Breakthrough Curve Analysis in Fixed-Bed Adsorption, Industrial \& Engineering Chemistry Research, vol. 55, 4734-4748 (2016). \href{https://doi.org/10.1021/acs.iecr.6b00516}{doi:10.1021/acs.iecr.6b00516}.

\bibitem{ward2022} Adam Ward, Ronny Pini, Integrated uncertainty quantification and sensitivity analysis of single-component dynamic column breakthrough experiments, Adsorption, vol. 28, 161-183 (2022). \href{https://doi.org/10.1007/s10450-022-00361-z}{doi:10.1007/s10450-022-00361-z}.

\bibitem{taylor2026} W. Taylor, J.D. Herman, V.L. Morales, Global sensitivity of laboratory-scale groundwater reactive transport simulations for sorption and degradation: Breakthrough dynamics across Péclet--Damköhler regimes, Journal of Contaminant Hydrology, vol. 283, 105043 (2026). \href{https://doi.org/10.1016/j.jconhyd.2026.105043}{doi:10.1016/j.jconhyd.2026.105043}.

\bibitem{miyabe2000} Kanji Miyabe, Georges Guiochon, Determination of the lumped mass transfer rate coefficient by frontal analysis, Journal of Chromatography A, vol. 890, 211-223 (2000). \href{https://doi.org/10.1016/s0021-9673(00)00537-9}{doi:10.1016/s0021-9673(00)00537-9}.

\bibitem{stanhope2014} S. Stanhope, J. E. Rubin, D. Swigon, Identifiability of Linear and Linear-in-Parameters Dynamical Systems from a Single Trajectory, SIAM Journal on Applied Dynamical Systems, vol. 13, 1792-1815 (2014). \href{https://doi.org/10.1137/130937913}{doi:10.1137/130937913}.

\bibitem{gallo2026} Matteo Gallo, Fabio Anselmi, Paolo Lazzari, The dynamics of discoverability: how trajectories and priors shape equation recovery, arXiv preprint arXiv:2607.18490v2 (2026). \href{https://arxiv.org/abs/2607.18490}{arXiv:2607.18490}.

\bibitem{delgado2025} Diemen Delgado-Cano, Erick Kracht, Urban Fasel, Benjamin Herrmann, SINDy on slow manifolds, arXiv preprint arXiv:2507.00747 (2025). \href{https://arxiv.org/abs/2507.00747}{arXiv:2507.00747}.

\bibitem{klishin2024} Andrei A. Klishin, Joseph Bakarji, J. Nathan Kutz, Krithika Manohar, Statistical mechanics of dynamical system identification, Physical Review Research, vol. 7, 033181 (2025). \href{https://doi.org/10.1103/4d98-tdlp}{doi:10.1103/4d98-tdlp}; \href{https://arxiv.org/abs/2403.01723}{arXiv:2403.01723}.

\bibitem{muhammed2025} Ismaila Muhammed, Dimitris M. Manias, Dimitris A. Goussis, Haralampos Hatzikirou, Data-driven identification of biological systems using multi-scale analysis, PLOS Computational Biology, vol. 21, e1013193 (2025). \href{https://doi.org/10.1371/journal.pcbi.1013193}{doi:10.1371/journal.pcbi.1013193}.

\bibitem{prabhu2025} Siddharth Prabhu, Nick Kosir, Mayuresh V. Kothare, Srinivas Rangarajan, Derivative-Free Domain-Informed Data-Driven Discovery of Sparse Kinetic Models, Industrial \& Engineering Chemistry Research, vol. 64, 2601-2615 (2025). \href{https://doi.org/10.1021/acs.iecr.4c02981}{doi:10.1021/acs.iecr.4c02981}.

\bibitem{huang2021} Juntao Huang, Yizhou Zhou, Wen-An Yong, Data-driven discovery of multiscale chemical reactions governed by the law of mass action, Journal of Computational Physics, vol. 448, 110743 (2022). \href{https://doi.org/10.1016/j.jcp.2021.110743}{doi:10.1016/j.jcp.2021.110743}; \href{https://arxiv.org/abs/2101.06589}{arXiv:2101.06589}.

\bibitem{delahunt2021} Charles B. Delahunt, J. Nathan Kutz, A toolkit for data-driven discovery of governing equations in high-noise regimes, IEEE Access, vol. 10, 31210-31234 (2022). \href{https://doi.org/10.1109/access.2022.3159335}{doi:10.1109/access.2022.3159335}; \href{https://arxiv.org/abs/2111.04870}{arXiv:2111.04870}.

\bibitem{kaptanoglu2023} Alan A. Kaptanoglu, Lanyue Zhang, Zachary G. Nicolaou, Urban Fasel, Steven L. Brunton, Benchmarking sparse system identification with low-dimensional chaos, Nonlinear Dynamics, vol. 111, 13143-13164 (2023). \href{https://doi.org/10.1007/s11071-023-08525-4}{doi:10.1007/s11071-023-08525-4}; \href{https://arxiv.org/abs/2302.10787}{arXiv:2302.10787}.

\bibitem{nogueira2022} Idelfonso B. R. Nogueira, Vinicius V. Santana, Ana M. Ribeiro, Alírio E. Rodrigues, Using scientific machine learning to develop universal differential equation for multicomponent adsorption separation systems, The Canadian Journal of Chemical Engineering, vol. 100, 2279-2290 (2022). \href{https://doi.org/10.1002/cjce.24495}{doi:10.1002/cjce.24495}.

\bibitem{frandsen2025} Jesper Frandsen, Vinicius V. Santana, Peter Jul-Rasmussen, Idelfonso B.R. Nogueira, Jakob K. Huusom, Krist V. Gernaey, Jens Abildskov, A systematic screening of neural network-based hybrid models of adsorption in chromatography processes, AIChE Journal, vol. 71, e70045 (2025). \href{https://doi.org/10.1002/aic.70045}{doi:10.1002/aic.70045}.

\bibitem{jessop2026} Arthur Jessop, Mohammed Alsubeihi, Ben Moseley, Ashwin Kumar Rajagopalan, Differentiable Hybrid Modelling for Learning and Optimising Chemical Transport Processes from Experimental Data, arXiv preprint arXiv:2609.04011 (2026). \href{https://arxiv.org/abs/2609.04011}{arXiv:2609.04011}.

\bibitem{xu2022} Xiao Xu, Marta D'Elia, Christian Glusa, John T. Foster, Machine-learning of nonlocal kernels for anomalous subsurface transport from breakthrough curves, arXiv preprint arXiv:2201.11146 (2022). \href{https://arxiv.org/abs/2201.11146}{arXiv:2201.11146}.

\bibitem{shukla2025} Jay Shukla, Xiaohui Qu, Zubin Darbari, Marija Iloska, J. Anibal Boscoboinik, Qin Wu, Discovering CO adsorption and desorption pathways from chemical reaction neural network modeling of transient kinetics spectroscopy, The Journal of Physical Chemistry Letters, vol. 16, 3562-3570 (2025). \href{https://doi.org/10.1021/acs.jpclett.5c00665}{doi:10.1021/acs.jpclett.5c00665}.

\bibitem{praditia2022} Timothy Praditia, Matthias Karlbauer, Sebastian Otte, Sergey Oladyshkin, Martin V. Butz, Wolfgang Nowak, Learning Groundwater Contaminant Diffusion-Sorption Processes With a Finite Volume Neural Network, Water Resources Research, vol. 58, e2022WR033149 (2022). \href{https://doi.org/10.1029/2022WR033149}{doi:10.1029/2022WR033149}.

\bibitem{subraveti2022} Sai Gokul Subraveti, Zukui Li, Vinay Prasad, Arvind Rajendran, Can a computer ``learn'' nonlinear chromatography?: Physics-based deep neural networks for simulation and optimization of chromatographic processes, Journal of Chromatography A, vol. 1672, 463037 (2022). \href{https://doi.org/10.1016/j.chroma.2022.463037}{doi:10.1016/j.chroma.2022.463037}.

\bibitem{scheurer2026} Stefania Scheurer, Riccardo Frenner, Tim Brünnette, Sergey Oladyshkin, Wolfgang Nowak, Efficient confidence interval computation for physics-aware machine learning of diffusion-sorption models, Frontiers in Water, vol. 8, 1813791 (2026). \href{https://doi.org/10.3389/frwa.2026.1813791}{doi:10.3389/frwa.2026.1813791}.

\bibitem{haghpanah2023} Reza Haghpanah, Danny Shade, Fitting Adsorption Isotherms with Symbolic Regression, Industrial \& Engineering Chemistry Research, vol. 62, 22141-22148 (2023). \href{https://doi.org/10.1021/acs.iecr.3c02900}{doi:10.1021/acs.iecr.3c02900}.

\bibitem{sharlin2024} Samiha Sharlin, Tyler R. Josephson, In Context Learning and Reasoning for Symbolic Regression with Large Language Models, arXiv preprint arXiv:2410.17448 (2024). \href{https://arxiv.org/abs/2410.17448}{arXiv:2410.17448}.


\bibitem{fox2023} Charles Fox, Neil Tran, Nikki Nacion, Samiha Sharlin, Tyler R. Josephson, Incorporating background knowledge in symbolic regression using a computer algebra system, Machine Learning: Science and Technology, vol. 5, 025057 (2024). \href{https://doi.org/10.1088/2632-2153/ad4a1e}{doi:10.1088/2632-2153/ad4a1e}; \href{https://arxiv.org/abs/2301.11919}{arXiv:2301.11919}.

\bibitem{servia2024} Miguel Ángel de Carvalho Servia, Ilya Orson Sandoval, King Kuok (Mimi) Hii, Klaus Hellgardt, Dongda Zhang, Ehecatl Antonio del Rio Chanona, The automated discovery of kinetic rate models -- methodological frameworks, Digital Discovery, vol. 3, 954-968 (2024). \href{https://doi.org/10.1039/d3dd00212h}{doi:10.1039/d3dd00212h}; \href{https://arxiv.org/abs/2301.11356}{arXiv:2301.11356}.

\bibitem{servia2025} Miguel Ángel de Carvalho Servia, Ilya Orson Sandoval, King Kuok, Hii, Klaus Hellgardt, Dongda Zhang, Ehecatl Antonio del Rio Chanona, Constraint-Guided Symbolic Regression for Data-Efficient Kinetic Model Discovery, arXiv preprint arXiv:2507.02730 (2025). \href{https://arxiv.org/abs/2507.02730}{arXiv:2507.02730}.

\bibitem{loman2025} Torkel E Loman, Ruth E Baker, Functional and parametric identifiability for universal differential equations applied to chemical reaction networks, arXiv preprint arXiv:2510.14140 (2025). \href{https://arxiv.org/abs/2510.14140}{arXiv:2510.14140}.

\bibitem{katsoulas2026} Konstantinos Katsoulas, Federico Galvanin, Luca Mazzei, Eva Sorensen, A Diagnostic Procedure for Identifying Isotherm Models in Liquid Chromatography, Industrial \& Engineering Chemistry Research, vol. 65, 1277-1292 (2026). \href{https://doi.org/10.1021/acs.iecr.5c03704}{doi:10.1021/acs.iecr.5c03704}.

\bibitem{abuelmaaty2024} Ahmed E. Abu EL-Maaty, Mahmoud A. Abdalla, Mohamed Essalhi, Mahmoud M. Abdelnaby, Morsi M. Mahmoud, Mohamed A. Habib, Mohamed Antar, Rached Ben-Mansour, Kinetics of Water Adsorption in Metal-Organic Framework(MOF-303) for Adsorption Cooling Application, Energy Conversion and Management: X, vol. 24, 100694 (2024). \href{https://doi.org/10.1016/j.ecmx.2024.100694}{doi:10.1016/j.ecmx.2024.100694}.

\bibitem{sakanaka2023} Yuta Sakanaka, Shotaro Hiraide, Iori Sugawara, Hajime Uematsu, Shogo Kawaguchi, Minoru T. Miyahara, Satoshi Watanabe, Generalised analytical method unravels framework-dependent kinetics of adsorption-induced structural transition in flexible metal--organic frameworks, Nature Communications, vol. 14, 6862 (2023). \href{https://doi.org/10.1038/s41467-023-42448-3}{doi:10.1038/s41467-023-42448-3}.

\bibitem{coyle2026} Reid A. Coyle, Shyam Chand Pal, Peter Walther, Saeun Park, Bin Feng, Zhiling Zheng, Sustainable metal--organic framework water harvesters in the artificial intelligence era, ACS Materials Letters, vol. 8, 2391-2400 (2026). \href{https://doi.org/10.1021/acsmaterialslett.6c00580}{doi:10.1021/acsmaterialslett.6c00580}; \href{https://arxiv.org/abs/2605.29179}{arXiv:2605.29179}.

\bibitem{raue2009} A. Raue, C. Kreutz, T. Maiwald, J. Bachmann, M. Schilling, U. Klingmüller, J. Timmer, Structural and practical identifiability analysis of partially observed dynamical models by exploiting the profile likelihood, Bioinformatics, vol. 25, 1923-1929 (2009). \href{https://doi.org/10.1093/bioinformatics/btp358}{doi:10.1093/bioinformatics/btp358}.

\bibitem{fung2026} Lloyd Fung, Overcoming error-in-variable or stiffness problem in SINDy-like data-driven model discovery, ACM Transactions on AI for Science, 3831701 (2026). \href{https://doi.org/10.1145/3831701}{doi:10.1145/3831701}.

\bibitem{bortz2023} David M. Bortz, Daniel A. Messenger, Vanja Dukic, Direct estimation of parameters in ODE models using WENDy: weak-form estimation of nonlinear dynamics, Bulletin of Mathematical Biology, vol. 85, 110 (2023). \href{https://doi.org/10.1007/s11538-023-01208-6}{doi:10.1007/s11538-023-01208-6}; \href{https://arxiv.org/abs/2302.13271}{arXiv:2302.13271}.

\bibitem{paper1} A.~Batayhi, The coefficient map, not the basis: what limits neural-operator and reduced-order surrogates on unseen adsorbents, companion manuscript (2026); code, results and pre-registrations archived at Zenodo, version v1.0.4, \href{https://doi.org/10.5281/zenodo.23312204}{doi:10.5281/zenodo.23312204}.

\bibitem{liu2024} Ziming Liu, Yixuan Wang, Sachin Vaidya, Fabian Ruehle, James Halverson, Marin Solja\v{c}i\'c, Thomas Y. Hou, Max Tegmark, KAN: Kolmogorov-Arnold Networks, ICLR 2025 (arXiv 2024). \href{https://arxiv.org/abs/2404.19756}{arXiv:2404.19756}.

\bibitem{slote2026} Kevin Slote, Jeremie Fish, Erik Bollt, KANDy: Kolmogorov-Arnold Networks and Dynamical System Discovery, arXiv preprint arXiv:2602.20413 (2026). \href{https://arxiv.org/abs/2602.20413}{arXiv:2602.20413}.

\bibitem{spotorno2026} Enzo Nicolas Spotorno, Josafat Leal Filho, Antonio Augusto Medeiros Frohlich, Empirical Stability Analysis of Kolmogorov-Arnold Networks in Hard-Constrained Recurrent Physics-Informed Discovery, arXiv preprint arXiv:2602.09988 (2026). \href{https://arxiv.org/abs/2602.09988}{arXiv:2602.09988}.
\bibitem{golub1980} Gene H. Golub, Charles F. Van Loan, An analysis of the total least squares problem, SIAM Journal on Numerical Analysis, vol. 17(6), 883-893 (1980). \href{https://doi.org/10.1137/0717073}{doi:10.1137/0717073}.

\bibitem{golub1987} G. H. Golub, Alan Hoffman, G. W. Stewart, A generalization of the Eckart-Young-Mirsky matrix approximation theorem, Linear Algebra and its Applications, vol. 88-89, 317-327 (1987). \href{https://doi.org/10.1016/0024-3795(87)90114-5}{doi:10.1016/0024-3795(87)90114-5}.
\bibitem{fajardo2023} Oscar Fajardo-Fontiveros, Ignasi Reichardt, Harry R. De Los R\'{i}os, Jordi Duch, Marta Sales-Pardo, Roger Guimer\`{a}, Fundamental limits to learning closed-form mathematical models from data, Nature Communications, vol. 14, 1043 (2023). \href{https://doi.org/10.1038/s41467-023-36657-z}{doi:10.1038/s41467-023-36657-z}.

\bibitem{messenger2024} Daniel A Messenger, David M Bortz, Asymptotic consistency of the WSINDy algorithm in the limit of continuum data, IMA Journal of Numerical Analysis, vol. 45, 3264--3312 (2025). \href{https://doi.org/10.1093/imanum/drae086}{doi:10.1093/imanum/drae086}.

\bibitem{hastings2024} Jon Hastings, Thomas Lassitter, Zhiling Zheng, Saumil Chheda, J. Ilja Siepmann, Laura Gagliardi, Omar M. Yaghi, T. Grant Glover, High-Temperature Water Adsorption Isotherms and Ambient Temperature Water Diffusion Rates on Water Harvesting Metal--Organic Frameworks, The Journal of Physical Chemistry C, vol. 128, 11328--11339 (2024). \href{https://doi.org/10.1021/acs.jpcc.4c01733}{doi:10.1021/acs.jpcc.4c01733}.

\bibitem{lassitter2024} Thomas Lassitter, Nikita Hanikel, Dennis J. Coyle, Mohammad I. Hossain, Bryce Lipinski, Michael O'Brien, David B. Hall, Jon Hastings, Juan Borja, Travis O'Neil, S. Ephraim Neumann, David R. Moore, Omar M. Yaghi, T. Grant Glover, Mass transfer in atmospheric water harvesting systems, Chemical Engineering Science, vol. 285, 119430 (2024). \href{https://doi.org/10.1016/j.ces.2023.119430}{doi:10.1016/j.ces.2023.119430}.

\bibitem{ma2025} Chunye Ma, Jiawang Li, Guang Wang, Zhigang Li, Wei Su, Yanguang Zhou, Microscopic Kinetics of Water Adsorption in Metal--Organic Frameworks, Nano Letters, vol. 25, 5778--5783 (2025). \href{https://doi.org/10.1021/acs.nanolett.5c00343}{doi:10.1021/acs.nanolett.5c00343}.

\bibitem{champneys2025} Max D. Champneys, Timothy J. Rogers, BINDy: Bayesian identification of nonlinear dynamics with reversible-jump Markov-chain Monte Carlo, Proceedings of the Royal Society A: Mathematical, Physical and Engineering Sciences, vol. 481, 20240620 (2025). \href{https://doi.org/10.1098/rspa.2024.0620}{doi:10.1098/rspa.2024.0620}.

\bibitem{ruan2025} Kai Ruan, Yilong Xu, Ze-Feng Gao, Yang Liu, Yike Guo, Ji-Rong Wen, Hao Sun, Discovering physical laws with parallel symbolic enumeration, Nature Computational Science, vol. 6, 53--66 (2025). \href{https://doi.org/10.1038/s43588-025-00904-8}{doi:10.1038/s43588-025-00904-8}.

\bibitem{buchner2026} David B\"{u}chner, Ronny Pini, Unary and Binary Dynamic Column Breakthrough Experiments with Carbon Dioxide and Nitrogen Imaged by X-ray Computed Tomography, Langmuir, vol. 42, 14716--14732 (2026). \href{https://doi.org/10.1021/acs.langmuir.6c00250}{doi:10.1021/acs.langmuir.6c00250}.

\bibitem{wang2022pyomo} Jialu Wang, Alexander W. Dowling, Pyomo.DOE: An open-source package for model-based design of experiments in Python, AIChE Journal, vol. 68, e17813 (2022). \href{https://doi.org/10.1002/aic.17813}{doi:10.1002/aic.17813}.

\bibitem{cornelio2023} Cristina Cornelio, Sanjeeb Dash, Vernon Austel, Tyler R. Josephson, Joao Goncalves, Kenneth L. Clarkson, Nimrod Megiddo, Bachir El Khadir, Lior Horesh, Combining data and theory for derivable scientific discovery with AI-Descartes, Nature Communications, vol. 14, 1777 (2023). \href{https://doi.org/10.1038/s41467-023-37236-y}{doi:10.1038/s41467-023-37236-y}.

\bibitem{shumaylov2025} Zakhar Shumaylov, Peter Zaika, Philipp Scholl, Gitta Kutyniok, Lior Horesh, Carola-Bibiane Sch\"{o}nlieb, When is a System Discoverable from Data? Discovery Requires Chaos, arXiv preprint arXiv:2511.08860 (2025). \href{https://arxiv.org/abs/2511.08860}{arXiv:2511.08860}.

\bibitem{rummel2025} Nic Rummel, Daniel A. Messenger, Stephen Becker, Vanja Dukic, David M. Bortz, WENDy for Nonlinear-in-Parameters ODEs, arXiv preprint arXiv:2502.08881 (2025). \href{https://arxiv.org/abs/2502.08881}{arXiv:2502.08881}.
\end{thebibliography}
\end{document}
